\section{Introduction}\label{sec:intro}

\subsection{The problem}

A matrix of $d\times d$ entries of $b=\lceil\log_2d\rceil$ bits, written row by row on the input tape of a
deterministic multitape Turing machine, occupies $L=d^2b$ cells. The machine must write the same matrix column
by column on its output tape. Can it do so in time $O(L)$?

On a random-access machine, or on a machine with a two-dimensional work tape, transposition is trivial, and a
Boolean circuit transposes with no gates at all. On one-dimensional tapes the problem is open. Recursive block
exchange gives time $O(L\log d)$. The best upper bound claimed is $O(L(\log d)^{1-\kappa})$ for a very small
$\kappa>0$, in an unrefereed preprint of OpenAI~\cite{OpenAI2026mult} whose transposition primitive mixes data by
XOR in sheared address frames. No superlinear lower bound is known; indeed none is known on machines with two or more work tapes for
off-line multiplication, sorting or the DFT either~\cite[\S1.4]{HvdH2025} (the superlinear lower bounds that
follow from~\cite{PPST1983} concern problems that are hard for nondeterministic linear time~\cite{Grandjean1988}).
The classical lower bounds for transposition hold for restricted models only: one work
tape~\cite{DMS1991,DietzfelbingerMaass1993}, data movement without coding~\cite{Stoss1973}, and reversal-bounded
conservative machines~\cite{Kirchherr1988}.

The question matters because of integer multiplication. Fast multiplication on Turing machines runs through
multidimensional discrete Fourier transforms~\cite{SchonhageStrassen1971,HvdH2021}, and on tapes such transforms
need transpositions to bring the data of each dimension together. Harvey and van der Hoeven proved a converse:
transposition reduces to multiplication, with no overhead for entries of $\lceil\log_2d\rceil$ bits~\cite[Cor.~6.9]{HvdH2025}. (Their own linear-time statement~\cite[Thm.~6.11]{HvdH2025}, in its simplest case, takes entries of $\lceil\log_2m\rceil$ bits for matrices of at most $m$ bits; up to linear time this is the same problem, \cref{prop:entry-width}.) So a linear-time multiplication machine would give a linear-time transposer, and a proof that transposition needs
superlinear time would be the first superlinear lower bound for off-line multiplication on multitape Turing
machines.

We write~(T) for the statement that a linear-time transposer exists. Its formal version,
\lean{Transpose.LinearTimeTransposable}, is fixed in a frozen Lean file (\cref{sec:model}), together with a checker
that accepts a proof of~(T) or of its negation.

\subsection{Results}

We do not decide~(T). We prove what we can about it, formally, and isolate one precise conjecture
whose proof would refute~(T).

\paragraph{Locality, counting and isoperimetry (\cref{sec:locality,sec:blocks,sec:traces}).} If a set of time steps
must output bits of the input that it never reads, those bits must enter through its entry data: the states at its
entry points and the work symbols it inherits (the Cut Lemma, \cref{lem:cut}). Cutting tapes into blocks of a row's
size and time into slots of the same length turns every run of linear length into a bounded-degree network on $O(d)$ slots, the
\emph{block network}, and the Cut Lemma into an isoperimetric inequality on it (\cref{cor:isoperimetry}). For machines
whose head motion depends on the data, inputs with the same block-level trace (a coarse record of head motion) behave like
an oblivious machine; for a linear-time machine some trace is shared by a $2^{-O(db)}$ fraction of the inputs (the popular trace, \cref{thm:popular}), and the
inequalities hold on it (\cref{lem:cut-trace,cor:isoperimetry-trace}).

\paragraph{Passes (\cref{sec:passes}).} \emph{Theorem~P.} An oblivious transposer whose run splits into $r$ passes, in
each of which every head moves monotonically, satisfies $(k+1)^r\ge1+kd/(2+\sigma_M/b)$, where $\sigma_M$ is the number
of bits of a state plus $k$ work symbols (\cref{thm:P}); \emph{Theorem~P$'$} gives $(k+1)^r\ge1+kd/(2+(\sigma_M+\gamma_T)/b)$ for every
machine halting within $T$ steps, where $\gamma_T=(k+4)\clog{(2T+1)}$ is $O(kb)$ when $T\le\poly(L)$
(\cref{thm:Pprime}). So, for $k\ge1$ and $T\le\poly(L)$, $\Omega(\log d)$ passes are necessary, whatever the
machine does with its data, and this remains true if the machine is correct only on a $2^{-O(db)}$ fraction of the inputs
(\cref{thm:P-on,thm:Pprime-on}).

\paragraph{Macroscopic motion (\cref{sec:mstar}).} \emph{Theorem~M*.} Rescale the work-head trajectories of a
linear-time transposer (for adaptive machines, along its popular trace classes) by $d$ in space (in blocks) and in
time (in slots). Along subsequences they converge to
$1$-Lipschitz paths, and every limit satisfies an isoperimetric inequality~$(\star)$: for every partition of macroscopic time into finitely many intervals, grouped into classes of measure at most $\theta$, the heads return to their positions across different classes with total multiplicity at least $1-2\theta$. Consequently the limits are not all monotone, and the heads
together waste at least one unit of range (\cref{thm:mstar,thm:mstar-adaptive,thm:excess}). The theorem holds for
non-uniform machines and for machines correct on a $2^{-o(L)}$ fraction of the inputs (\cref{thm:mstar-robust}), and
it is a statement about walks satisfying a cut inequality, not about machines (\cref{thm:mstar-walk}).

\paragraph{The Tape Coding Conjecture (\cref{sec:codes,sec:conjecture}).} The block-level results above use the machine only
through the flow of information along its block network. A run yields a \emph{causal network code} on the block
network of its time--space diagram that solves \emph{complete exchange}: every row must reach every column
(\cref{lem:class-code,thm:fast-codes}). On these networks almost every source--sink pair is at distance
$\Omega(\log d)$ while the total capacity is $O(L)$, so a solving code beats multicommodity flow by a factor
$\Omega(\log d)$ (\cref{thm:code-gap}). The \emph{Tape Coding Conjecture}
(\cref{conj:TCC}) asserts that such codes do not exist, even if they need to be correct only on a $2^{-O(d\log d)}$
fraction of the inputs; equivalently, coding beats routing on tape networks by $o(\log d)$ (\cref{thm:tcc-forms}). We
prove:
\begin{itemize}[leftmargin=2em]
\item the conjecture implies $\neg$(T); a proof of it yields the frozen checker's answer file directly
(\cref{thm:reduction}). The reduction is one-way: a refutation gives codes, not machines;
\item much less suffices: no codes at co-entropy $Cd$ for every $C$, for infinitely many $d$, on diagrams that write
their output at the end (\cref{thm:coent-linear,thm:io});
\item under the conjecture every machine fails at every large size, and no non-uniform machine with $2^{O(d\log d)}$
states is fast and correct on a $2^{-O(d\log d)}$ fraction of the inputs (\cref{thm:tcc-consequences});
\item every counterexample to the conjecture obeys Theorem~M* (\cref{cor:cex-star});
\item its plain form (codes correct on all inputs), which already refutes oblivious linear-time transposers
(\cref{thm:oblivious}), implies the robust form for linear and nearly linear codes, for affine and low-degree sets,
for codes with few error patterns, and for constant-size hints (\cref{sec:plainrobust}).
\end{itemize}

\paragraph{Barriers (\cref{sec:barriers}).} The same statement on general bounded-degree graphs would imply superlinear
circuit lower bounds for multiplication (\cref{prop:circuit-barrier}), and its bit-level analogue on tape networks
would separate oblivious from non-oblivious machines off-line (\cref{prop:rotation-barrier}). A proof must therefore use
the shape of tapes; and, since the plain form already excludes oblivious linear-time transposers
(\cref{thm:oblivious}), a proof would give a superlinear lower bound for oblivious transposition, which known
techniques do not provide.

\paragraph{XOR-frame recursions and constructions (\cref{sec:frames,sec:construction}).} In an abstract model of the recursion of the
OpenAI preprint (the correspondence is informal, \cref{sec:frames}), every run realizing a gadget makes at least as many frame changes as the gadget has wires (\cref{thm:frame-change}); by a translation from gadgets to runs, which we prove in full but have not formalized, every gadget has cost $s\ge W$ (\cref{prop:gadget-runs}), while linear time would need $s<W$; integer frames cost more. We record the conditions a construction must meet and the one architecture that we could not rule out; its open
cores are index-coding questions on bounded-degree expanders.

\subsection{How to read this paper: the formal proofs}

The proofs are the Lean proofs. Every numbered theorem, lemma, corollary, proposition and definition, except the four propositions marked as not formalized, names its Lean declarations and the file that contains them, in a line of the form
\begin{center}
{\small\textsf{Lean:} \texttt{declaration, \dots}\hfill[\texttt{file}]}
\end{center}
Declarations are in the namespace \texttt{Transpose.Tape} unless written with their full name, and files are relative to
\leanfile{lean/Transpose/}. The text after a statement, labelled \emph{Outline of the formal proof}, is a faithful and
straightforward guide to the formal proof, written for human readers; it is a convenience, not a proof. Where the
English statement simplifies (for instance by omitting harmless edge cases), the Lean statement is authoritative, and
the formal statements are often slightly stronger. The four results that are not formalized (\cref{prop:entry-width,prop:circuit-barrier,prop:rotation-barrier,prop:gadget-runs}) are marked as such and come with complete proofs. Results from the
literature are cited and not reproved.

The library depends only on Lean~4 and Mathlib (both v4.34.1). It contains no \texttt{sorry}, no unsafe or native code
evaluation, and no axioms beyond Lean's standard ones (\texttt{propext}, \texttt{Classical.choice}, \texttt{Quot.sound}).
\Cref{app:verify} explains how to build it and how to check these claims, and \cref{app:index} lists every declaration
cited in the paper.

\subsection{Related work}

\paragraph{Lower bounds on Turing machines.} Superlinear lower bounds are known for one-tape machines, starting with
Hennie~\cite{Hennie1965} and Maass~\cite{Maass1985}; for matrix transposition on one-tape off-line machines they are due
to Dietzfelbinger, Maass and Schnitger~\cite{DMS1991,DietzfelbingerMaass1993}. For multitape machines the strongest
general result is the separation of deterministic and nondeterministic linear time by Paul, Pippenger, Szemer\'edi and
Trotter~\cite{PPST1983}. Restricted models give more: Stoß proved $\Omega(L\log d)$ for permutations by data
movement~\cite{Stoss1973}, Kirchherr $\Omega(\log d)$ reversals for transposition on conservative
machines~\cite{Kirchherr1988}, and Pippenger and Fischer simulated multitape machines obliviously with a logarithmic
overhead~\cite{PippengerFischer1979}. Our Theorem~P is a reversal lower bound that allows arbitrary coding; lower bounds
for streams with few passes over external memory~\cite{GroheSchweikardt2005,GHS2009} need long or index-tagged items and
do not imply it.

\paragraph{External memory and network coding.} Aggarwal and Vitter determined the I/O complexity of sorting and
transposition for algorithms that move indivisible items~\cite{AggarwalVitter1988}. Whether coding can help is the
question behind the network coding conjecture of Li and Li~\cite{LiLi2004}, which would have implied lower bounds for
multiplication circuits~\cite{AFKL2019}, external-memory sorting~\cite{FHLS2019} and two-tape oblivious machines~\cite{AHJKL2006}.
Braverman and He~\cite{BravermanHe2026} and, concurrently, Zhang, Li and Li~\cite{ZhangLiLi2026} refuted it: in
undirected networks coding can beat multicommodity flow, by a factor $(\log n)^{\varepsilon}$ on networks of size $n$ after amplification~\cite{BGS2017}. Our conjecture is a version of the network coding
conjecture restricted to tape networks and to the densest demand, where a gain of order $\log d$ is at stake.

\paragraph{Multiplication.} Multiplication of $n$-bit integers in time $O(n\log n)$ on multitape machines is due to Harvey and van der
Hoeven~\cite{HvdH2021}, after Sch\"onhage and Strassen~\cite{SchonhageStrassen1971} and F\"urer~\cite{Furer2009}. The
OpenAI preprints claim multiplication below $n\log n$ and transposition below $L\log d$ on multitape machines~\cite{OpenAI2026mult} and exact complex linear circuits of size $o(n\log n)$ for the DFT of length $n$ for infinitely many lengths~\cite{OpenAI2026fourier};
we use them only as motivation, and \cref{sec:frames} shows the limits of their transposition technique.

\subsection{Organization}

\Cref{sec:model} fixes the model and the frozen statement. \Cref{sec:locality,sec:blocks} prove the cut lemmas,
\cref{sec:passes} Theorems~P and~P$'$, \cref{sec:traces} the extension to adaptive machines, and \cref{sec:mstar}
Theorem~M*. \Cref{sec:codes,sec:conjecture,sec:plainrobust} develop codes on tape diagrams, the Tape Coding Conjecture,
and the relation between its plain and robust forms. \Cref{sec:barriers} relates (T) to integer multiplication and explains why the conjecture is hard,
\cref{sec:frames} treats XOR-frame recursions, \cref{sec:construction} the construction side, and \cref{sec:open} lists
open problems. The appendices give the index of Lean declarations and instructions for building and verifying the
library.
